% ---------------------------------------------------------------------------
%  One-page activity handout -- Lab S2: predict, scan, check
%  Build: xelatex s2_activity_sheet.tex
% ---------------------------------------------------------------------------
\documentclass[10pt,a4paper]{article}
\usepackage{iftex}
\ifPDFTeX\usepackage[utf8]{inputenc}\usepackage[T1]{fontenc}\usepackage{lmodern}\else
\usepackage{fontspec}\usepackage{unicode-math}\setmainfont{Latin Modern Roman}\setsansfont{Latin Modern Sans}\setmathfont{Latin Modern Math}\fi
\usepackage[a4paper,margin=15mm]{geometry}
\usepackage{amsmath,booktabs,array,siunitx,graphicx,xcolor,enumitem,tikz,tcolorbox,microtype}
\tcbuselibrary{skins}
\graphicspath{{../figures/}}
% ---------------------------------------------------------------------------
%  s2figs.tex -- shared TikZ drawings for Lab S2 (Influence Lines)
%  Geometry in mm: A=0, B=250, C=440, D=690, E=880, F=1130.
%  \ILscale (cm per mm) is set by the including document before use.
% ---------------------------------------------------------------------------
\usetikzlibrary{arrows.meta,calc,patterns,decorations.pathreplacing}
\providecommand{\ILscale}{0.0115}
\colorlet{s2load}{orange!80!black}
\colorlet{s2reac}{blue!65!black}
\colorlet{s2hinge}{violet!80!black}
\colorlet{s2neg}{red!75!black}

% support symbols (anchored at beam centre line)
% offsets are given in cm so that the symbols keep their size whatever the x scale
\newcommand{\ILpin}[1]{%
  \draw[thick] (#1,-0.08) -- ++(-0.18cm,-0.3cm) -- ++(0.36cm,0) -- cycle;
  \draw[thick] ($(#1,-0.39)+(-0.26cm,0)$) -- ++(0.52cm,0);
  \foreach \t in {-0.22,-0.11,0,0.11,0.22} \draw ($(#1,-0.39)+(\t cm,0)$) -- ++(-0.08cm,-0.1cm);}
\newcommand{\ILroller}[1]{%
  \draw[thick] (#1,-0.08) -- ++(-0.18cm,-0.3cm) -- ++(0.36cm,0) -- cycle;
  \draw[thick] ($(#1,-0.45)+(-0.1cm,0)$) circle [radius=0.06cm] ($(#1,-0.45)+(0.1cm,0)$) circle [radius=0.06cm];
  \draw[thick] ($(#1,-0.52)+(-0.26cm,0)$) -- ++(0.52cm,0);
  \foreach \t in {-0.22,-0.11,0,0.11,0.22} \draw ($(#1,-0.52)+(\t cm,0)$) -- ++(-0.08cm,-0.1cm);}
\newcommand{\ILhingeat}[1]{\fill[white] (#1,0) circle [radius=0.11cm]; \draw[thick,s2hinge] (#1,0) circle [radius=0.11cm];}

% ---- the bridge model (elevation), with halts and span dimensions -----------
% optional argument: extra TikZ code drawn in the same picture (loads, arrows)
\newcommand{\ILbridge}[1][]{%
\begin{tikzpicture}[x=\ILscale cm,y=1cm,>=Stealth,font=\small]
  \pgfmathsetmacro{\xB}{250}\pgfmathsetmacro{\xC}{440}\pgfmathsetmacro{\xD}{690}\pgfmathsetmacro{\xE}{880}\pgfmathsetmacro{\xF}{1130}
  % beams
  \draw[line width=2.4pt] (0,0) -- (430,0);
  \draw[line width=2.4pt] (450,0) -- (680,0);
  \draw[line width=2.4pt] (700,0) -- (1130,0);
  \ILhingeat{440}\ILhingeat{690}
  \ILpin{0}\ILroller{250}\ILroller{880}\ILpin{1130}
  % halts
  \foreach \x/\l in {0/A,250/B,440/C,690/D,880/E,1130/F}
    \node[below,font=\small\bfseries] at (\x,-0.62) {\l};
  \foreach \x/\l in {100/1,200/2,300/3,400/4,500/5,600/6,700/7,800/8,900/9,1000/10,1100/11}
    {\draw (\x,-0.05) -- (\x,-0.22); \node[below,font=\scriptsize] at (\x,-0.22) {\l};}
  % dimensions
  \draw[|<->|] (0,-1.15) -- (250,-1.15) node[midway,below,font=\scriptsize] {250};
  \draw[|<->|] (250,-1.15) -- (440,-1.15) node[midway,below,font=\scriptsize] {190};
  \draw[|<->|] (440,-1.15) -- (690,-1.15) node[midway,below,font=\scriptsize] {250};
  \draw[|<->|] (690,-1.15) -- (880,-1.15) node[midway,below,font=\scriptsize] {190};
  \draw[|<->|] (880,-1.15) -- (1130,-1.15) node[midway,below,font=\scriptsize] {250};
  \node[font=\scriptsize,anchor=west] at (1140,-1.15) {mm};
  #1
\end{tikzpicture}}

% a downward load arrow at x (mm) with label
\newcommand{\ILload}[3][1.1]{\draw[->,very thick,s2load] (#2,#1) -- (#2,0.14) node[pos=0,above,font=\small\bfseries,s2load] {#3};}
% an upward reaction arrow at x with label (positive); negative drawn downward in red
\newcommand{\ILreaction}[3][0.9]{\draw[->,very thick,s2reac] (#2,-0.95-#1) -- (#2,-0.72) node[pos=0,below,font=\small,s2reac] {#3};}
\newcommand{\ILreactionneg}[3][0.9]{\draw[->,very thick,s2neg] (#2,-0.72) -- (#2,-0.95-#1) node[below,font=\small,s2neg] {#3};}

% ---- influence line plot: #1 = reaction name, #2 = list of x/y vertices -----
% #3 (optional) extra code; ordinate scale 1 unit = \ILyscale cm
\providecommand{\ILyscale}{1.0}
\newcommand{\ILilplot}[3][]{%
\begin{tikzpicture}[x=\ILscale cm,y=\ILyscale cm,>=Stealth,font=\small]
  \draw[gray!50,dashed] (250,-0.95) -- (250,1.95) (440,-0.95) -- (440,1.95) (690,-0.95) -- (690,1.95) (880,-0.95) -- (880,1.95);
  \draw[->] (0,0) -- (1180,0) node[right,font=\scriptsize] {$x$ (mm)};
  \foreach \x/\l in {0/A,250/B,440/C,690/D,880/E,1130/F} \node[below,font=\scriptsize\bfseries] at (\x,-0.02) {\l};
  \draw[very thick,s2reac] plot coordinates {#3};
  \node[anchor=south west,font=\small] at (0,1.5) {IL for #2};
  #1
\end{tikzpicture}}

\renewcommand{\ILscale}{0.0115}
\setlist{nosep,leftmargin=*}
\setlength{\parindent}{0pt}\setlength{\parskip}{4pt}
\pagestyle{empty}
\newcommand{\il}[1]{\eta_{#1}}
\newcommand{\pbox}[1]{\begin{tcolorbox}[enhanced,colback=white,colframe=black!50,boxrule=0.4pt,arc=1pt,left=4pt,right=4pt,top=2pt,bottom=2pt,height=#1]\end{tcolorbox}}
\begin{document}
\begin{minipage}[c]{0.76\textwidth}
{\LARGE\bfseries Predict -- scan -- check}\\[3pt]
{\large Lab S2 Influence Lines: in-class activity (15 minutes, in pairs)}\\[6pt]
Name(s): \underline{\hspace{6cm}}\quad Group: \underline{\hspace{2cm}}\\[4pt]
Write your prediction \emph{first}, in ink. Then scan the code (it opens the \emph{Explorer} tab of the Lab S2 course page), test it, and write what you found and why. A wrong prediction with a good explanation of what you learned earns full credit.
\end{minipage}\hfill
\begin{minipage}[c]{0.2\textwidth}\centering\includegraphics[width=\linewidth]{qr_s2_explorer}\\{\scriptsize Lab S2 page, Explorer tab}\end{minipage}

\begin{center}\ILbridge\end{center}

\begin{enumerate}[label=\textbf{\arabic*.},itemsep=3pt]
\item \textbf{Sign of $R_A$.} (Explorer: press \emph{Unit load} under \emph{Loads}.) A single mass stands at halt 4 (on the cantilever, 150 mm past B). Is the reaction at A upward, zero or downward? Prediction and reason:
\pbox{1.35cm}
Explorer result ($\il{A}$ at $x = 400$): \underline{\hspace{2cm}}\quad A negative reaction means: \underline{\hspace{4.8cm}}

\item \textbf{Largest $R_B$.} Where must a single mass stand for the largest reaction at B, and will $R_B$ be smaller or larger than the mass itself? Prediction:
\pbox{1.25cm}
Explorer: position \underline{\hspace{1.4cm}}, $\il{B} = $ \underline{\hspace{1.4cm}}. Why can a reaction exceed the load? \underline{\hspace{3.9cm}}

\item \textbf{The far side.} The mass is at halt 8 (between D and E). Predict $R_A$ and $R_B$, then check. What does the hinge at D transmit when the suspended span is unloaded?
\pbox{1.35cm}

\item \textbf{Compound load.} Press \emph{Test 2 pair} under \emph{Loads} (\SI{5}{kg} and \SI{2.5}{kg}, \SI{100}{mm} apart). Predict the position of the \SI{5}{kg} mass that gives the most negative $R_A$ and the value in kgf; then find it by trial.
\pbox{1.25cm}
Explorer: position \underline{\hspace{1.5cm}}, $R_A = $ \underline{\hspace{1.5cm}} kgf. Which mass stands on the vertex C, and why does that give the extreme? \underline{\hspace{4cm}}

\item \textbf{Distributed load.} Tick \emph{shade the areas} for $R_A$. Predict the sign of $R_A$ when a uniform load covers the whole model; then read the three areas and the net value.
\pbox{1.25cm}
Explorer: areas $+$\underline{\hspace{1.2cm}} mm and \underline{\hspace{1.2cm}} mm, net \underline{\hspace{1.2cm}} mm, so $R_A = $ \underline{\hspace{1.2cm}}$\,w$.
\end{enumerate}

\vspace{2pt}
\begin{tcolorbox}[enhanced,colback=orange!6,colframe=orange!75!black,boxrule=0.5pt,arc=1.5pt,left=5pt,right=5pt,top=3pt,bottom=3pt]
\textbf{In the laboratory:} after Test 1, open the \emph{Lab} tab of the same page and type your gauge-A readings into the Test 1 table. Note the largest absolute difference from theory and at which halt it occurs: \underline{\hspace{2cm}} at halt \underline{\hspace{1cm}}. Is it within the $\pm 0.03$ expected from a $\pm 5$ mm placement error and $\pm 1$ division? \underline{\hspace{3cm}}
\end{tcolorbox}
\end{document}
